% !TeX TS-program = xelatex

\documentclass[compress,aspectratio=169]{beamer}

\usepackage{presentationtemplate}
\usepackage[fontsize=\footnotesize, askip=3mm]{terminal}
\usepackage[fontsize=\footnotesize, linenosfontsize=\tiny, askip=3mm]{mylisting}
\usepackage{myitemize}

\title{Проектирование на C}

\begin{document}

\frame[plain]{\titlepage}

\begin{frame}

    \frametitle{Условия и принципы проектирования на C}

    \begin{myitemize}

        \item В C нет классов, наследования и исключений —
            но код должен быть поддерживаемым и масштабируемым.

        \item Проектирование на C собирается из простых механизмов:
            \texttt{struct}, указатели, указатели на функции, \texttt{enum}.

        \item Идиома — устойчивое сочетание этих механизмов,
            которое повторяется из проекта в проект.

    \end{myitemize}

\end{frame}

\begin{frame}

    \frametitle{План}

    \begin{myitemize}

        \item opaque pointer — сокрытие представления;

        \item handle — число вместо указателя;

        \item init/deinit — жизненный цикл объекта;

        \item возврат ошибок — коды и out-параметры;

        \item function table — поведение как данные;

        \item tagged union — аналог \texttt{std::variant} на C.

    \end{myitemize}

\end{frame}

\begin{frame}[fragile]

    \frametitle{Opaque pointer}

    \begin{columns}[T]

        \begin{column}{0.55\textwidth}

            \begin{myitemize}

                \item Потребитель видит только указатель \texttt{sec\_t*}
                    и объявления функций.

                \item Определение \texttt{struct} — только в файле .c
                    (реализации).

                \item Представление можно менять без пересборки клиентов.

            \end{myitemize}

        \end{column}

        \begin{column}{0.4\textwidth}

            \myinputlisting[minted language=c]
                {presentations/04-c-design/}
                {opaque.h}

        \end{column}

    \end{columns}

\end{frame}

\begin{frame}[fragile]

    \frametitle{Opaque pointer}

    \begin{columns}[T]

        \begin{column}{0.48\textwidth}

        Пример реализации библиотеки\footnotemark{}.

        \end{column}

        \begin{column}{0.48\textwidth}

            \myinputlisting[minted language=c,
                minted options={fontsize=\scriptsize, breaklines,
                    breakanywhere, linenos, escapeinside=||,
                    numbersep=2pt}]
                {presentations/04-c-design/}
                {opaque.c}

        \end{column}

    \end{columns}

    \footnotetext{%
        Объект хранится в \texttt{static}-переменной.
        Это намеренное упрощение.
    }

\end{frame}

\begin{frame}[fragile]

    \frametitle{Handle (дескриптор)}

    \begin{myitemize}

        \item Handle — это индекс объекта в таблице реализации:
            потребитель получает число, а не указатель.

        \item Развитие opaque pointer: представление по-прежнему
            скрыто, но вместо указателя — обычное целое.

        \item Плюсы: легко передавать, хранить и сравнивать;
            «висящих» указателей на стороне клиента не бывает.

        \item Реализация: пул объектов + поиск свободного места
            + невалидное значение.

    \end{myitemize}

\end{frame}

\begin{frame}[fragile]

    \frametitle{Handle (дескриптор)}

    Пример заголовочного файла библиотеки.

    \bigbreak

    \myinputlisting[minted language=c]
        {presentations/04-c-design/}
        {handle.h}

\end{frame}

\begin{frame}[fragile]

    \frametitle{Handle (дескриптор)}

    Пример реализации библиотеки.

    \bigbreak

    \myinputlisting[minted language=c,
        minted options={fontsize=\scriptsize, breaklines,
            breakanywhere, linenos, escapeinside=||, numbersep=2pt}]
        {presentations/04-c-design/}
        {handle.c}

\end{frame}

\begin{frame}[fragile]

    \frametitle{Init/deinit}

    \begin{myitemize}

        \item Пара функций: \texttt{init} создаёт и
            инициализирует объект, \texttt{deinit} освобождает
            ресурсы.

        \item \texttt{init} может провалиться — вернёт
            \texttt{NULL}; результат \texttt{init}
            всегда проверяем.

        \item \texttt{deinit} освобождает всё, что занял
            \texttt{init}: память, файлы, соединения.

        \item Аналог конструктора/деструктора в C: RAII нет,
            поэтому требуется ручной контроль.

    \end{myitemize}

\end{frame}

\begin{frame}[fragile]

    \frametitle{Init/deinit}

    Пример заголовочного файла библиотеки.

    \bigbreak

    \myinputlisting[minted language=c]
        {presentations/04-c-design/}
        {lifecycle.h}

\end{frame}

\begin{frame}[fragile]

    \frametitle{Init/deinit}

    Пример реализации библиотеки.

    \bigbreak

    \myinputlisting[minted language=c,
        minted options={fontsize=\scriptsize, breaklines,
            breakanywhere, linenos, escapeinside=||, numbersep=2pt}]
        {presentations/04-c-design/}
        {lifecycle.c}

\end{frame}

\begin{frame}[fragile]

    \frametitle{Возврат ошибок}

    \begin{myitemize}

        \item Код возврата: 0 = OK, > 0 — код ошибки;
            явно и самоочевидно.

        \item errno-стиль: вернуть \texttt{-1}/\texttt{NULL},
            код — в переменной \texttt{errno} (POSIX);
            слабость: переменная общая внутри потока —
            любой вызов библиотеки может её перезаписать.

    \end{myitemize}

\end{frame}

\begin{frame}[fragile]

    \frametitle{Возврат ошибок}

    \begin{columns}[T]

        \begin{column}{0.48\textwidth}

            \begin{mytitledinplacelisting}[minted language=c]{Пример с кодом возврата}
typedef enum {
    ERR_OK = 0,
    ERR_EMPTY,
} err_t;

err_t parse(int *out,
            char *s);

int main(int argc, char **v) {
    int i = 0;
    if (parse(&i, v[1]) != 0) {
        // process error...
    }
    return 0;
}
            \end{mytitledinplacelisting}

        \end{column}

        \begin{column}{0.48\textwidth}

            \begin{mytitledinplacelisting}[minted language=c]{Пример с errno}
typedef enum {
    ERR_OK = 0,
    ERR_EMPTY,
} err_t;

extern err_t err;

int parse(char *s);

int main(int argc, char **v) {
    int i = parse(v[1]);
    if (err) {
        // process error
        // and discard i...
    }
}
            \end{mytitledinplacelisting}

        \end{column}

    \end{columns}

\end{frame}

\begin{frame}[fragile]

    \frametitle{Function table}

    \begin{myitemize}

        \item Структура из указателей на функции — таблица операций
            («vtable» на C).

        \item Объект = данные + указатель на свою
            таблицу операций.

        \item Разные «классы» — разные таблицы
            с общими подписями.

    \end{myitemize}

\end{frame}

\begin{frame}[fragile]

    \frametitle{Function table}

    \myinputlisting[minted language=c]
        {presentations/04-c-design/}
        {vtable.c}

\end{frame}

\begin{frame}[fragile]

    \frametitle{Tagged union}

    \begin{myitemize}

        \item Tagged union позволяет создать sum type
            из теории типов\footnotemark{}.

        \item Тип кодируется при помощи константы-тега,
            а само значение хранится в объединении
            (\texttt{union}).

        \item Инвариант: тег и значение записываются вместе.

    \end{myitemize}

    \footnotetext{В product type данные хранятся одновременно (И),
        а в sum-type в конкретный момент времени хранится только один из
        возможных вариантов (ИЛИ).}

\end{frame}

\begin{frame}[fragile]

    \frametitle{Tagged union: пример}

    \myinputlisting[minted language=c]
        {presentations/04-c-design/}
        {tagged.c}

\end{frame}

\begin{frame}

    \frametitle{Что получили}

    \begin{myitemize}

        \item Инкапсуляция: opaque pointer → handle.

        \item Жизненный цикл: init/deinit.

        \item Ошибки: коды возврата, out-параметры, errno.

        \item Поведение: function table («vtable» на C).

        \item Данные: tagged union (аналог \texttt{std::variant} на C).

        \item Идиомы комбинируются: API библиотеки на C можно
            устроить «как класс».

    \end{myitemize}

\end{frame}

\end{document}
